\graphicspath{{conclusion/fig/}}

\chapter{Conclusion and Future Work}
\label{chap:conclusion}

\providecommand{\devname}{OpenStat}
\section{Conclusion}
\label{sec:conclusion}

This project set out to design, build and validate a low-cost, wireless potentiostat for the DPV detection of efavirenz and oseltamivir in water. Against the objectives of Section~\ref{sec:project_aim_obj_scope} (Table~\ref{tab:testing_summary}), the \devname{} applies the DPV waveform to within one DAC code (Objective~1); reads every current range to within \SI{0.44}{\percent} after calibration but overloads above \SI{720}{\micro\ampere} instead of \SI{750}{\micro\ampere} when in normal mode (Objective~2, partly met). The OpenStat agreed with dummy cells to within \SI{1.9}{\percent} and with the PalmSens4 within that instrument's own repeatability, apart from a \SI{24}{\milli\volt} offset (Objective~3, partly met). The device can be controlled physically or from a web browser (Objective~4) for an estimated \SIrange{18}{24}{\hour} per charge (Objective~5). It costs R2\,226, against US\$3\,340 for the PalmSens4 (Objective~6). The electrochemical tests were limited by the repeatability of the electrode, and the target drugs were not measured, so the detection limit was not measured; the noise floor suggests roughly \SIrange{1}{5}{\micro\gram\per\litre} for the electronics alone (Section~\ref{sec:randles_test}).

\section{Future Work}
\label{sec:future_work}

The \SI{24}{\milli\volt} offset should be investigated first, by repeating a ferricyanide scan with a \SI{25}{\milli\volt} pulse (Section~\ref{sec:fc_shift}) and on other current ranges. A ferricyanide concentration series on both instruments would then give the sensitivity and detection limit, and the noise floor should be measured with Wi-Fi on and off. In the firmware, starting each averaging window \SI{20}{\milli\second} before its edge would reduce the timing error (Section~\ref{sec:sw_limitations}), raising $C_f$ to \SI{63}{\pico\farad} would improve the phase margin with small-capacitance electrodes (Section~\ref{sec:stability}), and peak detection would let the device report peak height and potential. The hardware needs a grounded, shielded electrode cable (Section~\ref{sec:hw_ground}), the high-power mode of the potentiostat amplifier for the \SI{200}{\ohm} range (Section~\ref{sec:sw_ranges}), and a charger with power-path control so that it can run while charging. The device should then be tested with efavirenz and oseltamivir, using preconcentration or a modified electrode; machine learning could then be used to identify and quantify compounds whose peaks overlap.